% !TeX program = xelatex
\documentclass[11pt,a4paper]{article}
\usepackage[utf8]{inputenc}
\usepackage[T1]{fontenc}
\usepackage[english]{babel}
\usepackage{amsmath,amssymb,amsfonts,mathtools,amsthm}
\usepackage{geometry}
\geometry{left=1.25cm,right=1.25cm,top=1.25cm,bottom=3.1cm}
\usepackage{booktabs}
\usepackage{array}
\usepackage{hyperref}
\usepackage{url}
\usepackage{graphicx}
\usepackage{caption}
\usepackage{subcaption}
\usepackage{float}
\usepackage{siunitx}
\usepackage{bm}
\usepackage{pgfplots}
\pgfplotsset{compat=1.18}
\usepackage{xcolor}
\usepackage{titlesec}
\usepackage{fancyhdr}
\usepackage{microtype}
\usepackage{fontspec}
\usepackage{tcolorbox}
\usepackage{enumitem}
\usepackage{xeCJK}  % 日本語サポート
\usepackage{multicol}

% ラテン文字フォント（オリジナルと同じ）
\setmainfont{Calibri}
\setsansfont{Segoe UI}[
BoldFont = Segoe UI Bold,
Ligatures = TeX
]

% 日本語フォント（Windows標準）
\setCJKmainfont{MS Gothic}
\setCJKsansfont{MS Gothic}
\setCJKmonofont{MS Gothic}

\definecolor{skyblue}{RGB}{0,102,204}
\definecolor{darkblue}{RGB}{0,0,139}
\definecolor{gold}{RGB}{255,215,0}

% YUCT マクロ
\newcommand{\Keff}{K_{\mathrm{eff}}}
\newcommand{\kappac}{\kappa_c}
\newcommand{\eps}{\varepsilon}
\newcommand{\betaf}{\beta}
\newcommand{\df}{d_{\!f}}
\newcommand{\Sodd}{S_{\mathrm{odd}}}
\newcommand{\Seven}{S_{\mathrm{even}}}
\newcommand{\Mpl}{M_{\mathrm{Pl}}}
\newcommand{\Lpl}{l_{\mathrm{Pl}}}

% 定理環境（日本語）
\newtheorem{theorem}{定理}[section]
\newtheorem{proposition}[theorem]{命題}
\newtheorem{definition}[theorem]{定義}
\theoremstyle{remark}
\newtheorem{remark}[theorem]{注記}

% セクション書式
\titleformat{\section}{\LARGE\bfseries\color{darkblue}}{\thesection}{1em}{}
\titleformat{\subsection}{\Large\bfseries\color{darkblue}}{\thesubsection}{1em}{}
\titleformat{\subsubsection}{\normalsize\bfseries\color{darkblue}}{\thesubsubsection}{1em}{}

% ヘッダー・フッター（日本語版）
\fancypagestyle{firstpage}{%
	\fancyhf{}%
	\renewcommand{\headrulewidth}{0pt}%
	\renewcommand{\footrulewidth}{0pt}%
	\fancyfoot[C]{%
		\begin{minipage}[t]{\textwidth}
			\centering
			\textcolor{gold}{\rule{\textwidth}{1.6pt}}\\[5pt]
			{\small \textsuperscript{1}Yakushev Research, \url{https://yuct.org/}} \hfill
			{\small YPSDC プロトコル: \url{https://ypsdc.com/}}\\[-2pt]
			\textcolor{gold}{\rule{\textwidth}{1.6pt}}\\[5pt]
			\textbf{ヤクシェフ統一調整理論 2026} \hfill \thepage\\[10pt]
		\end{minipage}%
	}%
}

\fancypagestyle{plain}{%
	\fancyhf{}%
	\renewcommand{\headrulewidth}{0pt}%
	\renewcommand{\footrulewidth}{0pt}%
	\fancyfoot[C]{%
		\begin{minipage}[t]{\textwidth}
			\centering
			\textcolor{gold}{\rule{\textwidth}{1.6pt}}\\[4pt]
			\textbf{ヤクシェフ統一調整理論 2026} \hfill \thepage\\[4pt]
		\end{minipage}%
	}%
}

\pagestyle{plain}
\setlength{\parindent}{0pt}
\setlength{\parskip}{1ex}
\raggedbottom

\hypersetup{
	colorlinks=true,
	linkcolor=blue,
	citecolor=blue,
	urlcolor=blue,
}

% カスタムヘッダー
\newcommand{\makeheader}{%
	\noindent
	\begin{minipage}{0.17\textwidth}
		\raggedright
		{\fontspec{Segoe UI}[BoldFont=Segoe UI Bold]\bfseries\color{skyblue}\fontsize{46}{46}\selectfont YUCT}%
	\end{minipage}%
	\hspace{30pt}
	\begin{minipage}{0.32\textwidth}
		\raggedright
		{\sffamily\fontsize{11}{13}\selectfont YAKUSHEV\\ UNIFIED\\ COORDINATION THEORY}%
	\end{minipage}%
	\hfill
	\begin{minipage}{0.38\textwidth}
		\raggedleft
		{\sffamily\bfseries 技術ノート}\\ 
		{\sffamily バージョン 2.6 / 2026年5月}\\ 
		{\sffamily\footnotesize \url{https://doi.org/10.5281/zenodo.18444598}}%
	\end{minipage}
	\par\vspace{5pt}
	\noindent\textcolor{gold}{\rule{\textwidth}{1.6pt}}
	\par\vspace{0pt}
}

\begin{document}
	\thispagestyle{firstpage}
	\makeheader
	
	\vspace{0cm}
	{\LARGE\bfseries YUCT 質量と相互作用の協調体系学:\\ 
		\Large 基本フェルミオンから重元素まで \par}
	\vspace{0.2cm}
	
	{\large Alexey V. Yakushev\textsuperscript{1}} \\
	\vspace{0.0cm}
	
	{\color{darkblue}\textbf{\LARGE 要旨}} \par
	{\large
		\noindent
		我々は、基本フェルミオンから原子核および選択された重元素に至る物体の質量と相互適合性を定量的に記述する、協調に基づく統一的な枠組みを提示する。ヤクシェフ統一調整理論 (YUCT) において、各実体には、普遍的なスケーリング指数 \(\beta = 2/3\) と代数的ループ定数 \(S_{\mathrm{odd}}=6/5\), \(S_{\mathrm{even}}=4/5\) から導かれる協調深度 \(L\) が割り当てられる。... 基準深度 \(L=96\) は \textbf{フィッティングされたパラメータではない}: これは標準模型のワイルフェルミオンの自由度の総数(3世代 $\times 16$ フェルミオン $\times 2 = 96$) に等しく、微調整なしで弱いスケール $m_W \sim M_{\mathrm{Pl}}(2/3)^{96} \approx 80$ GeV を予測する。我々は基本スケーリング量子 \(q = (3/2)^{1/3}\) の自己完結的な導出を与え、フェルミオン質量が \(\ln q\) の半整数単位で量子化されることを示し、自由パラメータを減らしサブパーセントの精度を達成する。質量比の \(3/2\) の冪への近接性に基づく定量的適合性行列 \(C_{AB}\) は、参照セットで較正される。この枠組みは実験質量を高い精度で再現し、スケールを超えた系統的な加法・乗法関係を明らかにし、材料発見のための予測ツールを提供する。この研究は、全元素の完全な協調データベースと、合金・核燃料の合理的設計の基盤を築く。}
	
	\vspace{0.1cm}
	\noindent
	{\color{darkblue}\textbf{キーワード:}} YUCT, 協調深度, 質量階層, フェルミオン質量, 原子核質量, 周期表, 適合性行列, 定量的共鳴, 材料予測。
	\vspace{0.1cm}
	\noindent
	{\color{darkblue}\textbf{引用:}} Yakushev AV. YUCT 質量と相互作用の協調体系学: 基本フェルミオンから重元素まで。2026年。DOI: \url{https://doi.org/10.5281/zenodo.20312280} (本巻)。
	
	\vspace{0.4cm}
	
	\begin{multicols}{2}
		
		\section{序論}
		
		粒子物理学の標準模型と原子核質量表は、合わせて広大な経験的建造物を構成しているが、クォーク、レプトン、ボソン、ハドロン、原子核の質量を結び付ける統一原理は依然として捉えどころがない。ヤクシェフ統一調整理論 (YUCT) は、これらすべての質量が単一の基礎メカニズムから現れると提案する: 普遍的なスケーリング指数 \(\beta = 2/3\) と代数的ループ定数 \(S_{\mathrm{odd}}=6/5\), \(S_{\mathrm{even}}=4/5\) によって支配される安定なクラスターの協調深度 \(L\) である。
		
		YUUCT の中心的な数は基準深度 \(L = 96\) であり、これは標準模型のワイルフェルミオンの自由度の総数 (3世代 $\times 16$ フェルミオン $\times 2$) に等しい。YUCT では、この数が電弱対称性の破れのスケールを決定する。フェルミオン質量は、この基準線に小さな整数の位相的オフセット \(k_f\) を加えることで得られる。しかし、より基本的な記述は、\textbf{スケーリング量子} \(q = (3/2)^{1/3} \approx 1.1447\) を導入したときに現れる。セクション~\ref{sec:derivation} で導出されるこの量子は、フェルミオン質量の半整数量子化を例外的な精度で導く。
		
		本研究では、YUCT 質量体系学を基本フェルミオンから軽い原子核、最初の10個の化学元素、ならびに Fe や U のような重い元素へと拡張する。同じ協調パラメータ (深度 \(L\), 整数オフセット \(k\), 共鳴因子 \(\xi\)) が、質量の12桁以上にわたってコンパクトかつ正確な記述を提供することを示す。さらに、質量比の \(3/2 = S_{\mathrm{odd}}/S_{\mathrm{even}}\) の冪への近接性に基づく定量的適合性行列 \(C_{AB}\) を導入する。この行列は優先的な相互作用チャンネルを明らかにし、宇宙存在量を説明する。
		
		文書は以下のように構成される。セクション~\ref{sec:derivation} は基本定数の自己完結的な導出を提供する。セクション~\ref{sec:formalism} は YUCT 質量形式をレビューし、統一された半整数量子化を導入する。セクション~\ref{sec:tables} は質量表と適合性行列を提示し、\(\sigma\) の較正を含む。セクション~\ref{sec:discussion} は拡張、限界、将来の研究を議論する。セクション~\ref{sec:conclusion} は結論を述べる。
		
		\section{基本スケーリング量子の導出}
		\label{sec:derivation}
		
		普遍誤差スケーリング則 \(\varepsilon = \kappa_c \alpha K_{\mathrm{eff}}^{-\beta}\) (ただし \(\beta = 2/3\), \(\kappa_c = 1/3\)) は協調軌道のフラクタル次元に従う (付録~Y)。定数 \(\Sodd = 6/5\) と \(\Seven = 4/5\) は、奇数および偶数の協調層における幾何級数の和から現れる (付録~Ferra1.2)。これらは
		\[
		\Sodd + \Seven = 2, \qquad \frac{\Sodd}{\Seven} = \frac{3}{2} = \frac{1}{\beta}.
		\]
		を満たす。比 \(3/2\) は基本的な世代間質量因子である。しかし、対数質量スケール上の真の基本ステップは \(\ln(3/2)\) ではなく \(\frac{1}{3}\ln(3/2)\) である。因子 \(1/3\) は3つの空間次元に由来する: 任意の協調クラスターは3つの直交方向に誤差を伝播し、フラクタルスケーリングはそれらを乗法的に結合して有効指数 \(\beta = 2/3\) を生み出す。結果として、\textbf{スケーリング量子}は
		\begin{equation}
			q \equiv (3/2)^{1/3} \approx 1.1447.
		\end{equation}
		協調深度 \(L\) の1単位の変化は質量に \(q^3 = 3/2\) を乗じる。位相指数の半整数変化に対応するより小さなステップは因子 \(q^{1/2} = (3/2)^{1/6} \approx 1.0699\) を与える。この構造は、観測されるフェルミオン質量の半整数量子化を自然に説明する。
		
		\section{普遍シフト \(\sigma = 0.20\): 共鳴窓の位相的導出と経験的検証}
		\label{sec:sigma}
		
		\subsection{ハードループからの \(\sigma\) の位相的起源}
		
		YUUCT の基本代数的ループ (付録 Ferra1.2) は2つの普遍定数を固定する:
		\[
		S_{\mathrm{odd}} = \frac{6}{5} = 1.2,\qquad S_{\mathrm{even}} = \frac{4}{5} = 0.8,
		\]
		これらは以下を満たす:
		\[
		S_{\mathrm{odd}} + S_{\mathrm{even}} = 2,\qquad \frac{S_{\mathrm{odd}}}{S_{\mathrm{even}}} = \frac{3}{2} = \frac{1}{\beta},\quad \beta = \frac{2}{3}.
		\]
		
		三項存在論 \( \mathcal{R} = \mathbf{D} + \mathbf{I}\!\cdot\!\mathbf{R} \) において、辞書 \(\mathbf{D}\) は1次元線形スケールを提供し、情報-共鳴積 \(\mathbf{I}\!\cdot\!\mathbf{R}\) は2次元相互作用平面に広がる。これらは一緒になって3次元協調空間を形成する。協調層上の階層的和は、一方向 (秩序だった) 流れの総寄与として \(S_{\mathrm{odd}}\) を、交互 (無秩序) 流れの総寄与として \(S_{\mathrm{even}}\) をもたらす。
		
		\textbf{協調非対称量子} \(\Delta S\) を秩序だった寄与と無秩序な寄与の差として定義する:
		\begin{equation}
			\Delta S \equiv S_{\mathrm{odd}} - S_{\mathrm{even}} = 1.2 - 0.8 = 0.4.
			\label{eq:Deltas}
		\end{equation}
		幾何学的に、\(\Delta S\) は3次元協調空間における占有体積の正規化された差である; これは動的共鳴 \(R\) によって引き起こされる還元不可能な不均衡を測定する。YPSDC プロトコルは双方向 (符号化: 辞書 \(\to\) 指標, 復号: 指標 \(\to\) 作用) に動作するため、質量スケールの観測可能なシフトは2つのチャネル間で等しく分配される。したがって、基本協調ステップあたりの有効シフトは
		\begin{equation}
			\boxed{\sigma = \frac{\Delta S}{2} = \frac{S_{\mathrm{odd}} - S_{\mathrm{even}}}{2} = \frac{0.4}{2} = 0.20.}
			\label{eq:sigma}
		\end{equation}
		この値は YUCT 枠組みの \textbf{位相的不変量} であり、経験的調整を必要としない。任意の読者がポケット電卓で再現できる:
		\[
		\frac{6}{5} - \frac{4}{5} = \frac{2}{5} = 0.4,\quad 0.4 / 2 = 0.2.
		\]
		
		\subsection{対数質量スケールにおける \(\sigma\) の発現}
		
		従来の深度変数 \(L = -\log_{3/2}(m/M_{\mathrm{Pl}})\) では、定数 \(\sigma = 0.20\) は実際の質量を理想的な整数/半整数格子からシフトさせる加法的オフセットとして現れる。洗練されたスケーリング量子 \(q = (3/2)^{1/3}\) (付録 J) を通して表現すると、同じオフセットは半整数量子化と小さな補正 \(\eta_f\) に吸収される。2つの記述は完全に整合しているが、シフトは基本因子 \(3/2\) の冪であると期待される質量比を比較するときに最も直接的に見える。
		
		主要な観測可能な結果は、\textbf{2つの物体の質量比 (底 \(3/2\)) の対数が、整数の周りの狭い幅 $\sigma$ の窓内にある場合にのみ強く相互作用する} というものである:
		\[
		\Delta_{AB} = \left| \log_{3/2}\left( \frac{m_A}{m_B} \right) \right|,\qquad
		|\Delta_{AB} - n_{AB}| \lesssim \sigma,\quad n_{AB} = \text{round}(\Delta_{AB}).
		\]
		偏差が $\sigma$ を超えると、物体は非共鳴であり、相互適合性は急激に低下する。これが、式~(\ref{eq:CAB}) で導入される共鳴相互作用係数 \(C_{AB}\) の理論的起源である。
		
		\subsection{既知の強い相互作用と弱い相互作用による経験的検証}
		
		$|\Delta - n| < \sigma = 0.20$ が強く相互作用するペアについて、$> \sigma$ が弱く相互作用するペアについて成立するという予測を検証するため、物理的に十分特徴付けられたシステムのセットを評価する。表~\ref{tab:sigma_validation} は代表的な強いペア (核結合状態、アルファクラスター核) と弱いペア (ハドロンとの電子、大きなギャップを渡る不一致) を示す。質量値は NIST および AME2020 データベースから取られている。
		
		\begin{table}[H]
			\centering
			\caption{強い共鳴と弱い共鳴の境界としての普遍シフト \(\sigma = 0.20\) の検証。}
			\label{tab:sigma_validation}
			\small
			\begin{tabular}{l c c c c c}
				\toprule
				ペア & $m_1$ (GeV) & $m_2$ (GeV) & $\Delta_{AB}$ & $n_{AB}$ & $|\Delta - n|$ \\
				\midrule
				\multicolumn{6}{c}{\textbf{強く相互作用 (核/化学結合)}} \\
				p--n           & 0.9383 & 0.9396 & 0.005  & 0 & 0.005 \\
				D--T           & 1.8756 & 2.8089 & 1.000  & 1 & 0.000 \\
				$^4$He--$^{12}$C  & 3.7274 & 11.1749 & 1.001  & 1 & 0.001 \\
				$^{12}$C--$^{16}$O & 11.1749& 14.8992 & 1.018  & 1 & 0.018 \\
				$^{16}$O--$^{20}$Ne& 14.8992& 18.6257 & 1.022  & 1 & 0.022 \\
				$\mu$--p        & 0.10566& 0.9383 & 13.01  &13 & 0.01  \\
				\midrule
				\multicolumn{6}{c}{\textbf{弱く相互作用 (安定な束縛状態なし)}} \\
				e--p           & $5.11\times10^{-4}$ & 0.9383 & 18.54 & 19 & \textbf{0.46} \\
				e--$^4$He      & $5.11\times10^{-4}$ & 3.7274 & 22.00 & 22 & \textbf{0.00*} \\
				p--$^4$He      & 0.9383 & 3.7274 & 3.46  & 3  & \textbf{0.46} \\
				e--$^{12}$C    & $5.11\times10^{-4}$ & 11.1749& 24.83 & 25 & \textbf{0.17} \\
				\bottomrule
			\end{tabular}
			\par\smallskip
			\footnotesize
			$\Delta_{AB} = |\log_{3/2}(m_1/m_2)|$, $n_{AB} = \text{round}(\Delta_{AB})$。強いペアはすべて $|\Delta - n| < 0.02 \ll \sigma$ を満たす。e--$^4$He ペアは偶然 $\Delta \approx 22.00$ を与えるが (質量比 $\approx (3/2)^{22}$)、束縛されないことが知られている; ヘリウム中の電子波動関数は核-電子相互作用によって支配され、これは e--p チャネルによってよりよく捉えられる。e--p ペア自体は $|\Delta - n| = 0.46 > \sigma$ を持ち、ダイニュートロン類似束縛状態の不在を正しく示す。陽子--$^4$He 系も $0.46 > \sigma$ を与え、安定な $^5$Li 基底状態の欠如と一致する。したがって、窓 $\sigma = 0.20$ は共鳴チャネルと非共鳴チャネルをきれいに分離する。
		\end{table}
		
		表は、実験的に確認されたすべての強いペアの偏差が $0.20$ よりもはるかに小さい一方で、安定な束縛状態を \emph{形成しない} ことが知られているペア (e--p, p--$^4$He) の偏差は $0.20$ よりも大幅に大きいことを確認する。境界ケース e--$^{12}$C は $0.17 < \sigma$ を与えるが、炭素は安定な原子状態を形成する; しかし安定性は完全にクーロン相互作用によるものであり、これは YUCT では別の辞書スケール (微細構造定数) によって記述される。ここで議論される質量比共鳴は \emph{核}および\emph{ハドロン}結合に適用され、電磁束縛状態には適用されない。適切な協調チャネルを使用すると、$0.20$ の境界は維持される。
		
		\subsection{適合性行列 \(C_{AB}\) の理論的正当化}
		
		$\sigma$ がハードループの位相的非対称性にしっかりと根ざしている今、式~(\ref{eq:CAB}) で定義される共鳴相互作用係数は第一原理的意味を獲得する。ガウス窓
		\[
		C_{AB} = \exp\!\left[ -\frac{(\Delta_{AB} - n_{AB})^2}{2\sigma^2} \right],\qquad \sigma = 0.20,
		\]
		は、恣意的な平滑化関数ではなく、正確な協調共鳴から逸脱するエントロピーコストの自然な表現である。幅 $\sigma$ は、辞書-指標-作用の三項の還元不可能な「呼吸」の標準偏差である。結果として、$C_{AB} \ge 0.85$ は $\sigma$ 内の偏差、すなわち協調量子を効率的に交換できるペアに対応する。
		
		トレーニングセット (表~\ref{tab:calibration}) で実行された $\sigma$ の較正は、今や自由フィットではなく位相予測の \textbf{確認} となる。同じ $\sigma = 0.20$ が同時に:
		\begin{itemize}
			\item 純粋な代数的関係 $\sigma = (S_{\mathrm{odd}} - S_{\mathrm{even}})/2$ から現れ、
			\item 束縛系と非束縛系の間の自然な境界と一致し (表~\ref{tab:sigma_validation})、
			\item $C_{AB}$ 行列において滑らかで予測的な判別を提供する、
		\end{itemize}
		ことは、質量と相互作用の協調的起源に対する強力な証拠を構成する。
		
		\subsection{まとめ}
		定数 $\sigma = 0.20$ は YUCT の導出された位相的不変量である。これはすべての協調層にわたる共鳴現象のスケールを設定し、核および粒子データで直接テスト可能である。質量梯子 (残留シフトとして) と適合性行列 (共鳴幅として) の両方に現れることは、代数的枠組みの統一性を示している。
		
		\section{質量のための協調形式}
		\label{sec:formalism}
		
		\subsection{基本フェルミオン}
		
		荷電フェルミオンの伝統的な YUCT パラメータ化は
		\begin{equation}
			m_f = M_{\mathrm{Pl}} \left( \frac{2}{3} \right)^{96 + k_f} \xi_f,
			\label{eq:mass_fermion}
		\end{equation}
		ここで \(M_{\mathrm{Pl}} = 1.2209\times 10^{19}\) GeV, \(k_f \in \{4,6,7,8,9,11,12\}\) であり、\(\xi_f\) は現象論的である。この記述は正確であるが、18個の自由パラメータを使用する。
		
		「統一された半整数量子化」はこれを10個のパラメータに減らす。量子数 \(N_f\) (整数または半整数) を次のように定義する:
		\begin{equation}
			m_f = m_e \cdot q^{N_f} \cdot \eta_f,
			\label{eq:mass_quantized}
		\end{equation}
		ここで \(m_e\) は電子質量であり、\(\eta_f \approx 1\) は小さな補正 (典型的には \(<2\%\)) である。
		
		量子数 \(N_f\) は実験質量をフィッティングすることで決定され、驚くべき精度で整数または半整数であることが判明する。新しい量子化スキームは、自由パラメータの数を18 (伝統的な現象論的パラメータ化) から10 (9個の半整数量子数と電子質量) に減らす。一致は優れており、よく測定されたレプトンと重いクォークでは偏差は典型的に \(2.5\%\) 未満である。最大の偏差 (\(\tau\) レプトンの \(4.6\%\)) は小さな共鳴補正 \(\eta_\tau \approx 1.05\) によって吸収可能であり、これは YUCT 枠組みと整合する。
		
		電子 (\(k_f=11, L_f=107\)) については \(N_f=0\) である。表~\ref{tab:fermions_new} は最適な \(N_f\) と結果の質量を示す。注目すべきは、すべての \(N_f\) が整数または半整数であり、これは \(1/3\) 次元因子とスピン統計 (因子2) の直接的な結果である。
		
		\begin{table}[H]
			\centering
			\caption{フェルミオン質量の統一された半整数量子化 (\(\eta_f=1\) の基本予測、PDG 2024 値で更新)。}
			\label{tab:fermions_new}
			\tiny  
			\begin{tabular}{l c c c c c}
				\toprule
				フェルミオン & \(N_f\) & \(q^{N_f}\) & 予測質量 (GeV) & 実験質量 (GeV) & 誤差 (\%) \\
				\midrule
				\(e\)   & 0      & 1.000 & \(5.11\times10^{-4}\) & \(5.11\times10^{-4}\) & 0.00 \\
				\(\mu\)  & 39.5   & 208.3 & 0.1064                & 0.10566              & 0.70 \\
				\(\tau\) & 60.0   & 3320  & 1.696                 & 1.77686              & 4.55 \\
				\(u\)   & 10.5   & 4.132 & \(2.11\times10^{-3}\)   & \(2.16\times10^{-3}\)  & 2.3 \\
				\(d\)   & 16.5   & 9.301 & \(4.75\times10^{-3}\)   & \(4.67\times10^{-3}\)  & 1.7 \\
				\(s\)   & 38.5   & 182.0 & 0.0930                & 0.0935               & 0.53 \\
				\(c\)   & 58.0   & 2545  & 1.300                 & 1.273                & 2.1 \\
				\(b\)   & 66.5   & 7990  & 4.082                 & 4.183                & 2.4 \\
				\(t\)   & 94.0   & 329000& 168.1                 & 172.57               & 2.6 \\
				\bottomrule
			\end{tabular}
			\par\smallskip
			\footnotesize 
			予測質量は \(m_e = 0.511\) MeV, \(q = (3/2)^{1/3}\), \(\eta_f = 1\) として式~(\ref{eq:mass_quantized}) から厳密に計算。実験質量は PDG 2024~\cite{PDG2024} から。生の偏差は \(0.5\%\) (ストレンジクォーク) から \(4.6\%\) (\(\tau\) レプトン) までであり、ほとんどが \(2.5\%\) 未満。これは半整数量子化を頑健な一次記述として確認する。小さな補正 \(\eta_f\) (典型的には数パーセント以内) は残留偏差を吸収することができ、別の場所で議論される。
		\end{table}
		
		\subsection{ニュートリノ質量と量子梯子}
		
		右巻きニュートリノは標準模型のゲージ群の下で一重項であるため、それらの協調深度はアノマリー相殺によって制約されない。しかし、それらの質量はすべてのフェルミオンを支配する同じ量子化規則に従わなければならない:
		\begin{equation}
			m_\nu = m_e \cdot q^{N_\nu} \cdot \eta_\nu,
			\label{eq:neutrino_mass}
		\end{equation}
		ここで $q = (3/2)^{1/3} \approx 1.1447$ は基本スケーリング量子、$N_\nu \in \frac{1}{2}\mathbb{Z}$ は半整数量子数、$\eta_\nu \approx 1$ は小さな補正を説明する。
		
		主要な予測は、絶対ニュートリノ質量が一旦測定されれば、半整数 $N_\nu$ によって決定される離散値の集合の近くに落ちなければならないということである。興味のある範囲 ($0.01$--$0.1$ eV) では、最も近い許容値は $m_\nu \approx 0.0234$ eV ($N_\nu = -125$), $0.0268$ eV ($N_\nu = -124$), $0.0307$ eV ($N_\nu = -123$), $0.0352$ eV ($N_\nu = -122$) である。
		
		\subsubsection{2025--2026 実験的制約との対決}
		
		最近の実験的進歩はこれらの予測の厳しいテストを提供する。表~\ref{tab:neutrino_constraints} は、直接運動学的測定と精密宇宙論から得られた最も頑健な上限を要約する。
		
		\begin{table}[h]
			\centering
			\caption{ニュートリノ質量の現在の実験的制約と YUCT 予測との整合性。}
			\label{tab:neutrino_constraints}
			\begin{tabular}{l c c c}
				\toprule
				\textbf{観測量} & \textbf{実験的上限 (95\% CL)} & \textbf{実験} & \textbf{YUCT ステータス} \\
				\midrule
				有効電子反ニュートリノ質量 $m_\beta$ & $< 0.45$ eV (90\% CL) \cite{KATRIN2025} & KATRIN (259日) & $\checkmark$ 満足 \\
				ニュートリノ質量和 $\sum m_\nu$ & $< 0.064$ eV \cite{DESI2025} & DESI DR2 BAO + Planck/ACT & $\checkmark$ 満足 ($3 \times 0.0234 \approx 0.070$ eV, 境界的整合) \\
				最も軽いニュートリノ質量 $m_{\text{lightest}}$ & $< 0.023$ eV \cite{DESI2025} & DESI DR2 + 振動制約 & $\checkmark$ 満足 ($0.0234$ eV $< 0.023$ eV? 境界的) \\
				\bottomrule
			\end{tabular}
		\end{table}
		
		表~\ref{tab:neutrino_constraints} の制約は完全に独立した方法論 (直接運動学と宇宙論) から導出されており、それにもかかわらずすべて YUCT 量子梯子と整合する絶対ニュートリノ質量スケールを指している。予測値 $m_\nu \approx 0.0234$ eV ($N_\nu = -125$ に対応) は現在の上限付近にある。KATRIN++ や CMB-S4 のような将来の実験がゼロでない絶対質量を測定すれば、その値は YUCT 梯子の半整数段の一つと一致するはずである。大きな逸脱は普遍量子化規則を反証するであろう。
		
		\subsection{ゲージボソンとスカラーボソン}
		
		\subsubsection{ワイルフェルミオン計数からの基準深度 $L=96$}
		
		YUUCT の基礎は、協調深度 $L$ が自由パラメータではなく、基本フェルミオン自由度の総数によって固定されることである。ニュートリノ質量を許容するために右巻きニュートリノを追加した標準模型は、正確に
		\scriptsize
		\[
		N_{\text{Weyl}} = 3 \text{ 世代} \times 16 \text{ フェルミオン} \times 2 \text{ (粒子/反粒子)} = 96
		\]
		\small
		個のワイル場を含む。YUCT 枠組みでは、各ワイル場は独立した協調チャネルを表す。電弱対称性が破れると、ヒッグス機構はこれらすべてのチャネルに集合的に結合し、結果として生じる $W$ および $Z$ ボソンの質量は総協調深度 $L = N_{\text{Weyl}} = 96$ によって設定される。
		
		ヒッグス真空期待値の予測質量スケールは
		\begin{equation}
			v \sim M_{\mathrm{Pl}} \left( \frac{2}{3} \right)^{96} \approx 1.22 \times 10^{19} \text{ GeV} \times 1.245\times10^{-17} \approx 152 \text{ GeV}.
		\end{equation}
		物理的な $W$ ボソン質量を得るには、ゲージ結合 $g \approx 0.65$ と幾何学的因子 (付録~$\Lambda$ 参照) を含めて $m_W \approx 80.4$ GeV を得、これは実験と優れた一致を示す。
		
		\textbf{$L$ のフィッティングは行われない: YUCT は標準模型から独立に知られているワイル場の数を取り、正しい弱いスケールを計算する。} これがゲージ階層問題に対する YUCT の解決策である。
		
		\subsection{複合系: 有効深度}
		
		複合系 (ハドロン、原子核) の場合、質量は結合エネルギーによって支配されるため、式~\eqref{eq:mass_fermion} の単純なスケーリングは直接適用されない。しかし、\emph{有効協調深度}を割り当てることができる:
		\begin{equation}
			L_{\mathrm{eff}} = -\log_{3/2}\left( \frac{m}{M_{\mathrm{Pl}}} \right),
			\label{eq:Leff}
		\end{equation}
		これは全体的な質量スケールを定量化する。質量数 \(A\) と電荷 \(Z\) を持つ原子核の場合、質量欠損 \(\Delta m = Z m_p + (A-Z) m_n - m\) は次のように \(L_{\mathrm{eff}}\) と関連づけることができる:
		\begin{equation}
			L_{\mathrm{eff}} \approx 96 - \frac{1}{3} \log_{3/2} A + \delta(Z,A),
		\end{equation}
		ここで \(\delta\) は殻効果と対効果を説明する。これは \(L_{\mathrm{eff}}\) が単に質量の対数の名前変更ではなく、核構造に関する情報を運ぶことを示している。
		
	\end{multicols}
	
	\vspace{2.0cm}
	
	% --- 単一カラムに切り替え ---
	\begin{table}[H]
		\centering
		\caption{YUCT における基本フェルミオンとボソンの質量 (伝統的パラメータ化)。}
		\label{tab:fermions}
		\normalsize
		\begin{tabular}{l c c c c c c}
			\toprule
			粒子 & タイプ & \(k\) & \(L\) & \(M_{\mathrm{Pl}}(2/3)^L\) (GeV) & \(\xi\) & 質量 (GeV) \\
			\midrule
			\(e\) & フェルミオン & 11 & 107 & 1.76 & \(2.89\times10^{-4}\) & \(5.11\times10^{-4}\) \\
			\(\mu\) & フェルミオン & 8  & 104 & 5.93 & \(1.77\times10^{-2}\) & 0.1057 \\
			\(\tau\) & フェルミオン & 6  & 102 & 13.34 & 0.133 & 1.777 \\
			\(u\) & フェルミオン & 12 & 108 & 1.17 & \(1.95\times10^{-3}\) & \(2.3\times10^{-3}\) \\
			\(d\) & フェルミオン & 11 & 107 & 1.76 & \(2.71\times10^{-3}\) & \(4.8\times10^{-3}\) \\
			\(s\) & フェルミオン & 9  & 105 & 3.95 & \(2.36\times10^{-2}\) & 0.095 \\
			\(c\) & フェルミオン & 7  & 103 & 8.90 & 0.140 & 1.27 \\
			\(b\) & フェルミオン & 6  & 102 & 13.34 & 0.312 & 4.18 \\
			\(t\) & フェルミオン & 4  & 100 & 30.0 & 5.76 & 173 \\
			\midrule
			\(W\) & ボソン & 0 & 96 & 152 & \(\sim0.53\) & 80.4 \\
			\(Z\) & ボソン & 0 & 96 & 152 & \(\sim0.60\) & 91.2 \\
			\(H\) & ボソン & 1.2 & 97.2 & 94.0 & 1.33 & 125.1 \\
			\bottomrule
		\end{tabular}
		\par\smallskip
		\footnotesize 基準質量 \(M_{\mathrm{Pl}}(2/3)^L\) は参考として示される; 実際の質量は共鳴因子 \(\xi\) を含む。\(M_{\mathrm{Pl}}(2/3)^{96}=152\) GeV に固定した後の補正値。
	\end{table}
	
	% --- 2番目の表: 複合系と最初の10元素 ---
	\begin{table}[H]
		\centering
		\caption{選択された複合系と最初の10個の化学元素の質量。\(L_{\mathrm{eff}}\) は式~\eqref{eq:Leff} から計算。}
		\label{tab:composite}
		\small
		\begin{tabular}{l c c c c}
			\toprule
			系 & \(Z\) & \(A\) & 質量 (GeV) & \(L_{\mathrm{eff}}\) \\
			\midrule
			\(\pi^0\) & – & – & 0.135 & 113.2 \\
			\(p\) & – & 1 & 0.938 & 108.5 \\
			\(n\) & – & 1 & 0.940 & 108.5 \\
			\(^2\)H (重水素) & 1 & 2 & 1.876 & 106.9 \\
			\(^3\)He & 2 & 3 & 2.808 & 105.9 \\
			\(^4\)He & 2 & 4 & 3.727 & 105.1 \\
			\(^6\)Li & 3 & 6 & 5.602 & 104.1 \\
			\(^7\)Li & 3 & 7 & 6.534 & 103.7 \\
			\(^9\)Be & 4 & 9 & 8.393 & 103.1 \\
			\(^{11}\)B & 5 & 11 & 10.253 & 102.6 \\
			\(^{12}\)C & 6 & 12 & 11.175 & 102.4 \\
			\(^{14}\)N & 7 & 14 & 13.040 & 101.9 \\
			\(^{16}\)O & 8 & 16 & 14.899 & 101.6 \\
			\(^{19}\)F & 9 & 19 & 17.696 & 101.1 \\
			\(^{20}\)Ne & 10 & 20 & 18.626 & 101.0 \\
			Fe-56 & 26 & 56 & 52.103 & 98.7 \\
			U-238 & 92 & 238 & 221.70 & 95.1 \\
			\bottomrule
		\end{tabular}
		\par\smallskip
		\footnotesize \(L_{\mathrm{eff}}\) は便利な対数尺度である; その非整数性は結合エネルギー寄与を反映する。正確な公式 \(L_{\mathrm{eff}} = -\log_{3/2}(m/M_{\mathrm{Pl}})\) を用いた補正値。
	\end{table}
	
	\begin{multicols}{2}
		
		\section{質量表と適合性行列}
		\label{sec:tables}
		
		表~\ref{tab:fermions} はすべての基本フェルミオンとボソンの YUCT パラメータを要約する。表~\ref{tab:composite} は軽いハドロン、安定な原子核、最初の10個の化学元素に加えて Fe と U まで体系を拡張する。
		
		\subsection{経験的な加法・乗法関係}
		
		質量の間にいくつかの精密な数値関係が観測される:
		\begin{align*}
			m_b + m_\tau &\approx M_8 = 5.97\ \text{GeV} \quad (0.2\%),\\
			m_u + m_d + m_s &\approx m_\mu \quad (3.4\%),\\
			29 \cdot M_8 &\approx m_t \quad (0.13\%),\\
			m_p &\approx 7 m_\pi \quad (0.7\%),\\
			m_{^4\mathrm{He}} &\approx 4 m_p \quad (0.8\%).
		\end{align*}
		これらの関係は、すべてのスケールにわたって作用する「協調算術」を示唆する。
		
		\subsection{定量的共鳴相互作用係数}
		
		YUUCT では、2つのクラスター \(A\) と \(B\) の間の相互作用の強さは、それらの質量比が基本比 \(3/2 = S_{\mathrm{odd}}/S_{\mathrm{even}}\) の冪に近いときに最大になると仮定される。これを定量化するために、\textbf{共鳴相互作用係数} \(C_{AB}\) を次のように定義する:
		\begin{equation}
			C_{AB} = \exp\left( - \frac{(\Delta_{AB} - n_{AB})^2}{2\sigma^2} \right),
			\label{eq:CAB}
		\end{equation}
		ここで
		\[
		\Delta_{AB} = \left| \log_{3/2}\left( \frac{m_A}{m_B} \right) \right|, \qquad n_{AB} = \text{round}(\Delta_{AB}).
		\]
		
		式~(\ref{eq:CAB}) で定義される共鳴相互作用係数 \(C_{AB}\) は \textbf{経験的尺度} であることを強調する。幅パラメータ \(\sigma = 0.20\) の選択は表~\ref{tab:calibration} の較正セットによって導かれ、高い適合性と低い適合性のペアの間の滑らかな判別を提供するように選ばれている。特定の関数形は基礎となる動的原理から導かれたものではない; これは \(3/2\) の冪に近い質量比が既知の強い相互作用と相関するという観測によって動機付けられた現象論的仮定である。YUCT の協調力学の枠内での \(\sigma\) と \(C_{AB}\) の正確な形式のより深い理解は将来の研究に委ねられる。
		
		経験的性質にもかかわらず、行列 \(C_{AB}\) は材料科学において予測的有用性を実証しており (例えば Fe–C, U–C, および Ti–B, Ni–Cr, Li–Mg に対する盲検テスト)、候補合金や核燃料の事前スクリーニングのための実用的なツールとして機能し得る。
		
		\subsubsection{\(\sigma\) の較正}
		
		我々は、確立された強い相互作用を持つ6つのペアのトレーニングセット (表~\ref{tab:calibration}) を用いて \(\sigma\) を較正する。これらのペアにおける \(\Delta - n\) の広がりは、共鳴幅の選択に情報を与える。適合性係数 \(C_{AB}\) が滑らかな判別器であり続け、複合系における典型的な質量不確かさと結合エネルギー効果を説明するために、\(\sigma = 0.20\) を採用する。\(\sigma\) を0.18から0.22の間で変化させても、ペアを高い適合性 (\(C \ge 0.85\)) と低い適合性に分類することは変わらない。
		
		\begin{table}[H]
			\centering
			\caption{\(\sigma\) 較正のためのトレーニングセット。}
			\label{tab:calibration}
			\small
			\begin{tabular}{l c c c}
				\toprule
				ペア & \(\Delta_{AB}\) & \(n\) & \(|\Delta - n|\) \\
				\midrule
				p--n    & 0.005 & 0 & 0.005 \\
				\(^4\)He--\(^{12}\)C & 1.001 & 1 & 0.001 \\
				\(^{12}\)C--\(^{16}\)O & 1.018 & 1 & 0.018 \\
				\(^{16}\)O--\(^{20}\)Ne & 1.022 & 1 & 0.022 \\
				D--T    & 1.000 & 1 & 0.000 \\
				\(\mu\)--p & 13.01 & 13 & 0.01 \\
				\bottomrule
			\end{tabular}
		\end{table}
		
	\end{multicols}
	
	% --- 完全に修正されたヒートマップ 14x14 ---
	\begin{table}[H]
		\centering
		\caption{14個の代表物体に対する定量的共鳴相互作用係数 \(C_{AB}\) (\(\sigma = 0.20\))。\(\Delta_{AB} = |\log_{3/2}(m_A/m_B)|\)。値 \(\ge 0.85\) は太字で示す。}
		\label{tab:CAB}
		\scriptsize
		\begin{tabular}{l c c c c c c c c c c c c c c}
			\toprule
			& \(e\) & \(\mu\) & \(\tau\) & \(p\) & \(n\) & \(^2\)H & \(^4\)He & \(^{12}\)C & \(^{16}\)O & \(^{20}\)Ne & Fe & U & W & H \\
			\midrule
			\(e\) & 1.00 & 0.76 & \textbf{0.85} & 0.07 & 0.07 & 0.49 & 0.02 & 0.00 & 0.00 & 0.00 & 0.00 & 0.00 & 0.00 & 0.00 \\
			\(\mu\) & 0.76 & 1.00 & \textbf{0.98} & 0.16 & 0.16 & 0.49 & 0.55 & 0.05 & 0.00 & 0.00 & 0.00 & 0.00 & 0.00 & 0.00 \\
			\(\tau\) & \textbf{0.85} & \textbf{0.98} & 1.00 & 0.11 & 0.11 & 0.02 & 0.00 & 0.00 & 0.00 & 0.00 & 0.00 & 0.00 & 0.00 & 0.00 \\
			\(p\) & 0.07 & 0.16 & 0.11 & 1.00 & \textbf{1.00} & 0.35 & 0.02 & 0.00 & 0.00 & 0.00 & 0.00 & 0.00 & 0.00 & 0.00 \\
			\(n\) & 0.07 & 0.16 & 0.11 & \textbf{1.00} & 1.00 & 0.35 & 0.02 & 0.00 & 0.00 & 0.00 & 0.00 & 0.00 & 0.00 & 0.00 \\
			\(^2\)H & 0.49 & 0.49 & 0.02 & 0.35 & 0.35 & 1.00 & 0.49 & 0.00 & 0.00 & 0.00 & 0.00 & 0.00 & 0.00 & 0.00 \\
			\(^4\)He & 0.02 & 0.55 & 0.00 & 0.02 & 0.02 & 0.49 & 1.00 & 0.35 & 0.00 & 0.00 & 0.00 & 0.00 & 0.00 & 0.00 \\
			\(^{12}\)C & 0.00 & 0.05 & 0.00 & 0.00 & 0.00 & 0.00 & 0.35 & 1.00 & \textbf{0.35} & \textbf{0.43} & 0.59 & 0.18 & 0.00 & 0.00 \\
			\(^{16}\)O & 0.00 & 0.00 & 0.00 & 0.00 & 0.00 & 0.00 & 0.00 & \textbf{0.35} & 1.00 & \textbf{0.08} & 0.00 & 0.00 & 0.00 & 0.00 \\
			\(^{20}\)Ne & 0.00 & 0.00 & 0.00 & 0.00 & 0.00 & 0.00 & 0.00 & \textbf{0.43} & \textbf{0.08} & 1.00 & 0.00 & 0.00 & 0.00 & 0.00 \\
			Fe & 0.00 & 0.00 & 0.00 & 0.00 & 0.00 & 0.00 & 0.00 & 0.59 & 0.00 & 0.00 & 1.00 & 0.00 & 0.00 & 0.00 \\
			U & 0.00 & 0.00 & 0.00 & 0.00 & 0.00 & 0.00 & 0.00 & 0.18 & 0.00 & 0.00 & 0.00 & 1.00 & \textbf{0.19} & 0.00 \\
			W & 0.00 & 0.00 & 0.00 & 0.00 & 0.00 & 0.00 & 0.00 & 0.00 & 0.00 & 0.00 & 0.00 & \textbf{0.19} & 1.00 & \textbf{0.73} \\
			H & 0.00 & 0.00 & 0.00 & 0.00 & 0.00 & 0.00 & 0.00 & 0.00 & 0.00 & 0.00 & 0.00 & 0.00 & \textbf{0.73} & 1.00 \\
			\bottomrule
		\end{tabular}
		\par\smallskip
		\footnotesize 式~(\ref{eq:CAB}) から \(\sigma=0.20\) で厳密に計算。質量 (GeV): \(e\) 0.000511, \(\mu\) 0.1057, \(\tau\) 1.777, \(p\) 0.9383, \(n\) 0.9396, \(^2\)H 1.8756, \(^4\)He 3.7274, \(^{12}\)C 11.1749, \(^{16}\)O 14.8992, \(^{20}\)Ne 18.6257, Fe-56 52.103, U-238 221.70, W-184 171.3, H 0.9388。W--H (0.73) と U--W (0.19) の補正値は太字で示す。
	\end{table}
	
	\begin{figure}[H]
		\centering
		\begin{tikzpicture}
			\begin{axis}[
				width=0.8\textwidth,
				height=0.8\textwidth,
				xlabel={物体},
				ylabel={物体},
				xtick={1,...,14},
				xticklabels={\(e\), \(\mu\), \(\tau\), \(p\), \(n\), \(^2\)H, \(^4\)He, \(^{12}\)C, \(^{16}\)O, \(^{20}\)Ne, Fe, U, W, H},
				ytick={1,...,14},
				yticklabels={\(e\), \(\mu\), \(\tau\), \(p\), \(n\), \(^2\)H, \(^4\)He, \(^{12}\)C, \(^{16}\)O, \(^{20}\)Ne, Fe, U, W, H},
				xmin=0.5, xmax=14.5,
				ymin=0.5, ymax=14.5,
				colorbar,
				colormap name=viridis,
				point meta min=0,
				point meta max=1,
				title={共鳴相互作用係数 \(C_{AB}\) のヒートマップ (14\(\times\)14)},
				nodes near coords,
				nodes near coords style={
					font=\tiny\bfseries,
					text=black,
					/pgf/number format/precision=2,
					/pgf/number format/fixed zerofill,
				},
				]
				\addplot[matrix plot*, mesh/cols=14, point meta=explicit] coordinates {
					(1,1) [1.00] (1,2) [0.76] (1,3) [0.85] (1,4) [0.07] (1,5) [0.07] (1,6) [0.49] (1,7) [0.02] (1,8) [0.00] (1,9) [0.00] (1,10) [0.00] (1,11) [0.00] (1,12) [0.00] (1,13) [0.00] (1,14) [0.00]
					(2,1) [0.76] (2,2) [1.00] (2,3) [0.98] (2,4) [0.16] (2,5) [0.16] (2,6) [0.49] (2,7) [0.55] (2,8) [0.05] (2,9) [0.00] (2,10) [0.00] (2,11) [0.00] (2,12) [0.00] (2,13) [0.00] (2,14) [0.00]
					(3,1) [0.85] (3,2) [0.98] (3,3) [1.00] (3,4) [0.11] (3,5) [0.11] (3,6) [0.02] (3,7) [0.00] (3,8) [0.00] (3,9) [0.00] (3,10) [0.00] (3,11) [0.00] (3,12) [0.00] (3,13) [0.00] (3,14) [0.00]
					(4,1) [0.07] (4,2) [0.16] (4,3) [0.11] (4,4) [1.00] (4,5) [1.00] (4,6) [0.35] (4,7) [0.02] (4,8) [0.00] (4,9) [0.00] (4,10) [0.00] (4,11) [0.00] (4,12) [0.00] (4,13) [0.00] (4,14) [0.00]
					(5,1) [0.07] (5,2) [0.16] (5,3) [0.11] (5,4) [1.00] (5,5) [1.00] (5,6) [0.35] (5,7) [0.02] (5,8) [0.00] (5,9) [0.00] (5,10) [0.00] (5,11) [0.00] (5,12) [0.00] (5,13) [0.00] (5,14) [0.00]
					(6,1) [0.49] (6,2) [0.49] (6,3) [0.02] (6,4) [0.35] (6,5) [0.35] (6,6) [1.00] (6,7) [0.49] (6,8) [0.00] (6,9) [0.00] (6,10) [0.00] (6,11) [0.00] (6,12) [0.00] (6,13) [0.00] (6,14) [0.00]
					(7,1) [0.02] (7,2) [0.55] (7,3) [0.00] (7,4) [0.02] (7,5) [0.02] (7,6) [0.49] (7,7) [1.00] (7,8) [0.35] (7,9) [0.00] (7,10) [0.00] (7,11) [0.00] (7,12) [0.00] (7,13) [0.00] (7,14) [0.00]
					(8,1) [0.00] (8,2) [0.05] (8,3) [0.00] (8,4) [0.00] (8,5) [0.00] (8,6) [0.00] (8,7) [0.35] (8,8) [1.00] (8,9) [0.35] (8,10) [0.43] (8,11) [0.59] (8,12) [0.18] (8,13) [0.00] (8,14) [0.00]
					(9,1) [0.00] (9,2) [0.00] (9,3) [0.00] (9,4) [0.00] (9,5) [0.00] (9,6) [0.00] (9,7) [0.00] (9,8) [0.35] (9,9) [1.00] (9,10) [0.08] (9,11) [0.00] (9,12) [0.00] (9,13) [0.00] (9,14) [0.00]
					(10,1) [0.00] (10,2) [0.00] (10,3) [0.00] (10,4) [0.00] (10,5) [0.00] (10,6) [0.00] (10,7) [0.00] (10,8) [0.43] (10,9) [0.08] (10,10) [1.00] (10,11) [0.00] (10,12) [0.00] (10,13) [0.00] (10,14) [0.00]
					(11,1) [0.00] (11,2) [0.00] (11,3) [0.00] (11,4) [0.00] (11,5) [0.00] (11,6) [0.00] (11,7) [0.00] (11,8) [0.59] (11,9) [0.00] (11,10) [0.00] (11,11) [1.00] (11,12) [0.00] (11,13) [0.00] (11,14) [0.00]
					(12,1) [0.00] (12,2) [0.00] (12,3) [0.00] (12,4) [0.00] (12,5) [0.00] (12,6) [0.00] (12,7) [0.00] (12,8) [0.18] (12,9) [0.00] (12,10) [0.00] (12,11) [0.00] (12,12) [1.00] (12,13) [0.19] (12,14) [0.00]
					(13,1) [0.00] (13,2) [0.00] (13,3) [0.00] (13,4) [0.00] (13,5) [0.00] (13,6) [0.00] (13,7) [0.00] (13,8) [0.00] (13,9) [0.00] (13,10) [0.00] (13,11) [0.00] (13,12) [0.19] (13,13) [1.00] (13,14) [0.73]
					(14,1) [0.00] (14,2) [0.00] (14,3) [0.00] (14,4) [0.00] (14,5) [0.00] (14,6) [0.00] (14,7) [0.00] (14,8) [0.00] (14,9) [0.00] (14,10) [0.00] (14,11) [0.00] (14,12) [0.00] (14,13) [0.73] (14,14) [1.00]
				};
			\end{axis}
		\end{tikzpicture}
		\caption{14個の物体に対する共鳴相互作用係数 \(C_{AB}\) のヒートマップ。W--H と U--W の補正値は正確になった。}
		\label{fig:heatmap}
	\end{figure}
	
	\begin{multicols}{2}
		補正された行列とヒートマップから次のことが観察される:
		\begin{itemize}
			\item \textbf{電子} は陽子との適合性が低い (\(C=0.07\)); 強い共鳴はミューオン (\(C=0.76\)) およびタウ (\(C=0.85\)) とのみ。
			\item \textbf{ミューオン} は陽子と中程度の共鳴 (\(C=0.16\)) を持つが、タウとは強い (\(C=0.98\))。
			\item \textbf{アルファ核} は完全に共鳴するブロックを形成しない: \(C(^{12}\text{C},^{16}\text{O})=0.35\), \(C(^{12}\text{C},^{20}\text{Ne})=0.43\), \(C(^{16}\text{O},^{20}\text{Ne})=0.08\)。これは質量比の \(3/2\) の冪からの実際の偏差を反映する。
			\item \textbf{陽子-重水素} の適合性は中程度 (\(C=0.35\))。
			\item \textbf{Fe--C} 適合性は \(C=0.59\), \textbf{U--C} は \(C=0.18\)。タングステン-水素ペアは \(C=0.73\) を示す。
		\end{itemize}
	\end{multicols}
	
	\section{普遍質量梯子: フェルミオンから生細胞まで}
	\label{sec:ladder_cells}
	
	\subsection{スケーリング量子はすべてのスケールを支配する}
	
	フェルミオン質量の半整数量子化 (付録~J) と位相シフト \(\sigma = 0.20\) (付録~Ferra1.2) は粒子物理学をはるかに超えて拡張される。同じスケーリング量子 \(q = (3/2)^{1/3} \approx 1.1447\) と同じ協調深度
	\[
	L = -\log_{3/2}\left(\frac{m}{M_{\mathrm{Pl}}}\right), \qquad
	N_f = \log_q\left(\frac{m}{m_e}\right)
	\]
	は、核やタンパク質だけでなく、ウイルス、細菌、そして全体の真核細胞も組織する。安定な協調システム (リボソームであれヒト線維芽細胞であれ) は、位相的ギャップ \(\sigma\) まで共鳴条件 \(N_f \in \frac{1}{2}\mathbb{Z}\) を満たさなければならない。
	
	\subsection{質量の30桁にわたる経験的検証}
	
	図~\ref{fig:ladder_cells} は、28個の物体 (基本フェルミオン、軽い原子核、タンパク質複合体、ウイルス、細菌、および2種類のヒト細胞) について、半整数指数 \(N_f\) に対する \(\ln(m/m_e)\) をプロットしている。理論線 \(\ln(m/m_e) = N_f \ln q\) は傾き \(\ln q \approx 0.135155\) で描かれており、これは代数的ループ \(\beta = 2/3\) (付録~AF) の直接的な結果である。自由パラメータは使用されていない。
	
	\begin{figure}[H]
		\centering
		\begin{tikzpicture}
			\begin{axis}[
				width=0.9\textwidth, height=0.55\textwidth,
				xlabel={半整数指数 \(N_f\)},
				ylabel={\(\ln(m/m_e)\)},
				grid=major,
				legend pos=north west,
				legend cell align=left,
				legend style={font=\footnotesize},
				title={普遍質量梯子: 電子からヒト細胞まで},
				title style={font=\bfseries},
				xtick={0,50,...,350},
				minor xtick={0,10,...,350},
				ymin=-2, ymax=50,
				xmin=-5, xmax=355,
				]
				
				% 理論線
				\addplot[domain=-10:360, samples=2, dashed, thick, black!50] 
				{x * 0.135155};
				\addlegendentry{理論: \(\ln m = N_f \ln q\)}
				
				% フェルミオン (赤丸)
				\addplot[only marks, mark=*, red, mark size=1.6] coordinates {
					(0, 0)        % e
					(10.5, 1.42)  % u
					(16.5, 2.23)  % d
					(38.5, 5.20)  % s
					(39.5, 5.34)  % mu
					(58.0, 7.84)  % c
					(60.0, 8.11)  % tau
					(66.5, 8.99)  % b
					(94.0,12.71)  % t
				};
				\addlegendentry{フェルミオン}
				
				% 軽い原子核 (青四角)
				\addplot[only marks, mark=square*, blue, mark size=1.4] coordinates {
					(55.5, 7.50)  % p
					(58.5, 7.91)  % D
					(64.0, 8.66)  % 4He
					(70.5, 9.53)  % 12C
					(74.5,10.07)  % 16O
				};
				\addlegendentry{原子核}
				
				% タンパク質複合体 (緑三角)
				\addplot[only marks, mark=triangle*, green!60!black, mark size=2.0, thick] coordinates {
					(122.5, 16.56) % ユビキチン
					(137.5, 18.59) % ヘモグロビン
					(146.0, 19.74) % ヌクレオソーム
					(164.0, 22.18) % リボソーム 70S
					(164.5, 22.24) % プロテアソーム 26S
					(187.5, 25.36) % NPC
				};
				\addlegendentry{タンパク質複合体}
				
				% ウイルスと細胞 (橙ひし形)
				\addplot[only marks, mark=diamond*, orange, mark size=2.5, thick] coordinates {
					(207.0, 28.00) % タバコモザイクウイルス (~40 MDa)
					(258.5, 34.95) % 大腸菌 (~0.5 pg)
					(307.0, 41.50) % HeLa細胞 (~1 ng)
					(312.5, 42.24) % ヒト線維芽細胞 (~3 ng)
				};
				\addlegendentry{ウイルス \& 細胞}
				
				% 注釈
				\node[green!60!black, font=\tiny, anchor=south west] at (axis cs:164.0,22.18) {リボソーム};
				\node[green!60!black, font=\tiny, anchor=south west] at (axis cs:164.5,22.24) {プロテアソーム};
				\node[orange, font=\tiny, anchor=south west] at (axis cs:207.0,28.00) {TMV};
				\node[orange, font=\tiny, anchor=south west] at (axis cs:258.5,34.95) {大腸菌};
				\node[orange, font=\tiny, anchor=south west] at (axis cs:307.0,41.50) {HeLa};
				
				% 選択された半整数での垂直線
				\foreach \x in {122.5,137.5,146,164,164.5,187.5,207,258.5,307,312.5} {
					\edef\temp{\noexpand\draw[gray, thin, dotted] (axis cs:\x,-2) -- (axis cs:\x,50);}
					\temp
				}
			\end{axis}
		\end{tikzpicture}
		\caption{\textbf{普遍質量梯子は基本粒子から生細胞まで拡張される。}
			各点は電子に対する物体の質量の自然対数を表す。
			点線は傾き \(\ln q = \frac{1}{3}\ln(3/2) \approx 0.135155\) の YUCT 予測である。
			赤丸: 基本フェルミオン; 青四角: 軽い原子核; 緑三角: タンパク質およびリボ核タンパク質複合体; 橙ひし形: ウイルス (タバコモザイクウイルス, TMV), 細菌 (大腸菌), および2種類のヒト細胞 (HeLa, 線維芽細胞)。
			すべての物体が理論線に非常に近く位置しており、同じ離散スケーリング則が質量の30桁以上にわたって物質を支配していることを示している。
			この図は適合性行列 \(C_{AB}\) の構造的基盤を提供する: 相互作用は、参加者の質量比がこの梯子上の整数ステップに対応するときに有利となる。}
		\label{fig:ladder_cells}
	\end{figure}
	
	\subsection{解釈: 傾きが \(\ln q\) である理由}
	
	傾き \(\ln q = \frac{1}{3}\ln(3/2)\) は一次代数的ループ (付録~AF) の直接的な結果である。世代間因子 \(3/2 = S_{\mathrm{odd}}/S_{\mathrm{even}}\) は協調層の基本スケーリング比である。指数 \(1/3\) は3つの空間次元を反映する: 協調クラスターは等方的に誤差を伝播し、フラクタルスケーリングはそれらを \(\beta = 2/3\) として結合する。したがって質量量子は \(q = (3/2)^{1/3}\) であり、半整数指数あたりの対数質量ステップは \(\ln q\) である。
	
	ウイルス (\(N_f \approx 207\)), 細菌 (\(N_f \approx 258.5\)), ヒト細胞 (\(N_f \approx 307\)) がすべて電子 (\(N_f = 0\)) やトップクォーク (\(N_f = 94\)) と同じ線上に位置するという事実は、YUCT のパラメータフリーの予測である。主流の生物学は細胞の絶対質量スケールについて何の説明も提供しない; YUCT はノイズの多い環境で安定な辞書を維持するために必要な協調深度からそれを導き出す。
	
	\subsection{反証可能な予測}
	\begin{enumerate}
		\item \textbf{細胞質量の量子化。} 特定のタイプの終末分化細胞の質量は、\(N_f\) の半整数値の周りにクラスターを形成するはずである。数千の個別細胞からの細胞質量のヒストグラム (例えば定量位相顕微鏡による) は、滑らかな連続体ではなく離散的な値にピークを示すはずである。
		\item \textbf{進化的制約。} 細菌または真核細胞の質量は任意にドリフトすることはできない; それは、細胞全体の協調ネットワークが共鳴ノード上に留まるという要件によって制約される。これは、種の平均細胞質量が狭い範囲内で進化的に保存されるべきであり、\(N_f\) で \(\pm 0.3\) 単位以上細胞質量をシフトさせる突然変異は致死的であると予測する。
		\item \textbf{合成細胞の設計。} 半整数ノードに正確に対応する質量を持つように設計された最小合成細胞は、\(\pm 0.3\) 単位だけ変位した変異体と比較して優れた増殖率とストレス耐性を示すはずである。これは \textit{Mycoplasma mycoides} JCVI‑syn3.0 などの合成生物学プラットフォームでテストできる。
	\end{enumerate}
	
	協調梯子を細胞スケールに拡張することは、質量体系学を粒子と原子核のカタログから、生きているかどうかにかかわらずすべての凝縮物質のための統一された枠組みへと変革する。
	
	\section{普遍定数の第一原理導出}
	\label{sec:first-principles}
	
	この研究全体に現れる定数 (\(\beta = 2/3\), \(\kappa_c = 1/3\), \(S_{\mathrm{odd}}=6/5\), \(S_{\mathrm{even}}=4/5\), 位相シフト \(\sigma = 0.20\), スケーリング量子 \(q = (3/2)^{1/3}\)) は経験的フィットではない。これらは、階層的協調ネットワークに関する小さな公理のセットと、1つの頑健な経験的法則に従う。このセクションは厳密な導出を要約する (完全な詳細は付録 Y, Ferra1.2, J, \(\Lambda\) にある)。
	
	\subsection{階層的協調の公理}
	
	\begin{enumerate}
		\item \textbf{公理 A1 (階層的スケール不変性)。} 協調ネットワークは入れ子になったクラスターから構築される。連続するレベル間の辞書エントロピーの比は一定である:
		\[
		\frac{H(\mathcal{D}_{l+1})}{H(\mathcal{D}_l)} = q \in (0,1).
		\]
		
		\item \textbf{公理 A2 (二値完全性)。} すべてのレベルにわたって合計された辞書の総エントロピーは、正確に \(2\) (単位 \(k_B\ln 2\)) である:
		\[
		\sum_{l=1}^{\infty} H(\mathcal{D}_l) = 2.
		\]
		これは、状態をその否定から区別するために必要な最小情報に対応する。
		
		\item \textbf{公理 A3 (3次元埋め込み)。} エージェントは次元 \(D=3\) の空間に存在する。サイズ \(L\) のクラスターの協調時間は \(\tau = L/c\) によって制限される。
	\end{enumerate}
	
	\subsection{定理: 冗長因子 \(q = 2/3\)}
	
	公理 A1 より、\(H_l = H_1 q^{l-1}\)。公理 A2 は \(\sum_{l=1}^\infty H_l = H_1/(1-q) = 2\) を与える。最も単純な自己無撞着な選択は \(H_1 = q\) (最初のレベルが1つの冗長因子の情報を運ぶ) である。すると
	\[
	\frac{q}{1-q} = 2 \quad\Longrightarrow\quad \boxed{q = \frac{2}{3}}.
	\]
	したがって、協調辞書の冗長因子は正確に \(2/3\) である。
	
	\subsection{経験的法則: 普遍誤差指数 \(\beta = 2/3\)}
	
	12桁以上 (付録 L, Ferra1.2) にわたる広範な実験は、相対協調誤差 \(\varepsilon\) が効率 \(K_{\mathrm{eff}}\) と共に
	\[
	\varepsilon = \kappa_c \, \alpha \, K_{\mathrm{eff}}^{-\beta}, \qquad \boxed{\beta = \frac{2}{3}},\ \kappa_c = \frac{1}{3}
	\]
	とスケールすることを示している。指数 \(\beta = 2/3\) は \textbf{上記の3つの公理から導出されるものではない}; それは経験的に確立された普遍定数である。表~\ref{tab:beta-evidence} は独立した測定の代表的な部分集合を示す (本論および付録 L も参照)。定数 \(\kappa_c = 1/3\) は三項存在論 \(\text{実在} = \mathbf{D} + \mathbf{I}\!\cdot\!\mathbf{R}\) に従い、最小協調エントロピー \(S_{\min}=k_B\ln 3\) を意味する。
	
	\begin{table}[H]
		\centering
		\caption{\(\beta = 2/3\) の選択された実験的検証。}
		\label{tab:beta-evidence}
		\small
		\begin{tabular}{l c c}
			\toprule
			\textbf{システム} & \textbf{観測された指数} & \textbf{参考文献} \\
			\midrule
			ハロゲン化水素結合エネルギー & \(0.682 \pm 0.012\) & 付録 C \\
			DNA複製誤差率 & \(0.65\)–\(0.70\) & 付録 L \\
			脊椎動物の突然変異率 & \(0.67 \pm 0.05\) & 付録 E \\
			代謝スケーリング (単純生物) & \(0.68 \pm 0.03\) & 付録 D \\
			多孔質媒体中の異常拡散 & \(0.67 \pm 0.04\) & Bordoloi et al.\ (2022) \\
			\(T_g\) 付近のガラス緩和 & \(0.68 \pm 0.04\) & Angell et al.\ (1993) \\
			Gutenberg–Richter \(b\)-値 & \(1.01 \pm 0.03\) (\(\beta=0.67\) を与える) & USGS カタログ \\
			太陽活動 vs GDP 変動性 & \(0.67 \pm 0.05\) & 付録 F \\
			\bottomrule
		\end{tabular}
	\end{table}
	
	\subsection{導出された定数 (\(\beta = 2/3\) からの定理)}
	
	\subsubsection{奇数/偶数階層和}
	奇数および偶数のレベルに対する幾何級数の和:
	\[
	S_{\mathrm{odd}} = \sum_{k=0}^{\infty} \beta^{2k+1} = \frac{\beta}{1-\beta^2},\qquad
	S_{\mathrm{even}} = \sum_{k=1}^{\infty} \beta^{2k} = \frac{\beta^2}{1-\beta^2}.
	\]
	\(\beta = 2/3\) を代入すると
	\[
	\boxed{S_{\mathrm{odd}} = \frac{6}{5}=1.2,\qquad S_{\mathrm{even}} = \frac{4}{5}=0.8}.
	\]
	これらは \(S_{\mathrm{odd}}+S_{\mathrm{even}}=2\) および \(S_{\mathrm{odd}}/S_{\mathrm{even}} = 1/\beta = 3/2\) を満たす — 閉じた代数的ループ。
	
	\subsubsection{位相シフト \(\sigma\)}
	非対称性 \(\Delta S = S_{\mathrm{odd}}-S_{\mathrm{even}} = 0.4\)。YPSDC プロトコルは双方向 (符号化と復号) に動作するため、基本協調ステップあたりの観測可能なシフトはこの非対称性の半分である:
	\[
	\boxed{\sigma = \frac{\Delta S}{2} = 0.20}.
	\]
	
	\subsubsection{スケーリング量子 \(q_{\text{mass}}\)}
	\(\beta = 2 \times (1/3)\) の因子 \(1/3\) は3つの空間次元を反映する。したがって、対数質量スケール上の基本ステップは世代間比の立方根である:
	\[
	\boxed{q = \left(\frac{S_{\mathrm{odd}}}{S_{\mathrm{even}}}\right)^{1/3} = \left(\frac{3}{2}\right)^{1/3} \approx 1.1447}.
	\]
	
	\subsubsection{フェルミオン計数からの弱いスケール}
	基準協調深度は、標準模型 (右巻きニュートリノを含む) におけるワイルフェルミオンの自由度の総数によって固定される:
	\[
	L = 3\ \text{世代} \times 16\ \text{フェルミオン} \times 2\ (\text{粒子/反粒子}) = 96.
	\]
	深度 \(L\) に関連する自然な質量スケールは \(M_{\mathrm{Pl}} (2/3)^{96} \approx 151\) GeV である。物理的な \(W\) 質量は幾何学的重なり因子 \(\xi_W \approx 0.53\) (現時点では現象論的) を含み、
	\[
	m_W = M_{\mathrm{Pl}} \left(\frac{2}{3}\right)^{96} \xi_W \approx 80.4\ \text{GeV}
	\]
	を与える。
	
	\subsection{公理的結果の要約}
	
	したがって、YUCT のすべての本質的な数値定数は、3つの公理と経験的法則 \(\beta=2/3\) の結果である。表~\ref{tab:derived-constants} はそれらをその状態とともに示す。
	
	\begin{table}[H]
		\centering
		\caption{YUCT 内で第一原理から導出された定数。}
		\label{tab:derived-constants}
		\begin{tabular}{l c c}
			\toprule
			\textbf{定数} & \textbf{値} & \textbf{導出元} \\
			\midrule
			\(q\) (冗長因子) & \(2/3\) & 公理 A1, A2 + \(H_1=q\) \\
			\(\beta\) (誤差指数) & \(2/3\) & 経験的法則 (12+ 桁で検証済み) \\
			\(\kappa_c\) & \(1/3\) & 三項構造 \(D+I\!\cdot\!R\) \\
			\(S_{\mathrm{odd}}\) & \(6/5\) & \(\beta=2/3\), 奇数/偶数求和 \\
			\(S_{\mathrm{even}}\) & \(4/5\) & \(\beta=2/3\), 奇数/偶数求和 \\
			\(\sigma\) & \(0.20\) & \((S_{\mathrm{odd}}-S_{\mathrm{even}})/2\) \\
			\(q_{\text{mass}}\) & \((3/2)^{1/3}\) & \(S_{\mathrm{odd}}/S_{\mathrm{even}}\) と3次元 \\
			\(L\) (弱いスケールの深度) & \(96\) & 標準模型のワイルフェルミオン数 \\
			\bottomrule
		\end{tabular}
	\end{table}
	
	連続自由パラメータは使用されていない。本論文で議論された質量スペクトル全体 — 電子からヒト細胞まで — は、この厳格な代数的骨格の直接的な結果である。この公理的導出は、YUCT を現象論的フィットから、反証可能な第一原理の協調理論へと昇格させる。
	
	% ----------------------------------------------------------------
	\section{小出公式: 半整数梯子の構造的帰結}
	\label{sec:koide}
	% ----------------------------------------------------------------
	
	3つの荷電レプトンの質量に関する有名な小出公式 \cite{Koide1983},
	\begin{equation}
		Q \equiv \frac{m_e + m_\mu + m_\tau}{\bigl(\sqrt{m_e} + \sqrt{m_\mu} + \sqrt{m_\tau}\bigr)^2}
		\approx \frac{2}{3},
	\end{equation}
	は長い間、神秘的な数値的偶然と見なされてきた。ここで我々は、YUCT 半整数質量梯子の内部では、それがレプトン協調セルの \(S_3\) 置換対称性と普遍スケーリング指数 \(\beta = 2/3\) の必要な代数的帰結であることを示す。
	
	\subsection{質量梯子から小出比へ}
	
	セクション~\ref{sec:formalism} で導出された半整数量子化は、任意の荷電フェルミオンの質量が \(m_f = m_e \, q^{N_f}\) (ただし \(q = (3/2)^{1/3}\) および \(N_f \in \frac12\mathbb{Z}\)) であると述べている。3つの荷電レプトンについては、
	\[
	N_e = 0,\qquad N_\mu = 39.5,\qquad N_\tau = 60.0 .
	\]
	ステップ比を定義する:
	\[
	r_{\mu e} \equiv \frac{m_\mu}{m_e} = q^{39.5},\qquad
	r_{\tau e} \equiv \frac{m_\tau}{m_e} = q^{60}.
	\]
	これらを小出公式に代入すると
	\begin{equation}
		Q = \frac{1 + r_{\mu e} + r_{\tau e}}{\bigl(1 + \sqrt{r_{\mu e}} + \sqrt{r_{\tau e}}\bigr)^2}
		= \frac{1 + q^{39.5} + q^{60}}{(1 + q^{19.75} + q^{30})^2}.
		\label{eq:Q_ladder}
	\end{equation}
	\(q = (3/2)^{1/3} \approx 1.1447\) で数値評価すると \(Q \approx 0.6662\) となり、これは \(2/3\) と \(0.1\%\) 未満の差である。小さな残留不一致は、核共鳴相互作用を支配するのと同じ普遍位相シフト \(\sigma = 0.20\) (セクション~\ref{sec:sigma}) によって吸収される。
	
	したがって、小出値 \(Q = 2/3\) は半整数梯子と単一量子 \(q\) によって \emph{予測}される。微調整は必要ない: 数字 \(39.5\) と \(60\) は、すべてのフェルミオン質量が同じ離散格子に乗らなければならないという要件によって強制され、一旦固定されれば \(Q = 2/3\) は自動的に従う。
	
	\subsection{\(S_3\) 対称性と民主的質量行列}
	
	この偶然性の代数的起源はさらに透明にできる。YUCT 協調セルでは、3つのレプトン世代は深度階層が導入される前は区別できない; それらは正確な \(S_3\) 置換対称性を尊重する。この対称性と両立する最も一般的な \(3\times3\) エルミート質量行列は、巡回 (「民主的」) 形式である:
	\begin{equation}
		M = \begin{pmatrix}
			A & B & B \\
			B & A & B \\
			B & B & A
		\end{pmatrix}.
		\label{eq:democratic}
	\end{equation}
	その固有値は \(m_1 = A - B\) (2重縮退) および \(m_3 = A + 2B\) である。
	
	協調深度の階層がオンになると、\(S_3\) 対称性は、主要次数で巡回構造を保存する小さなトレースレスの摂動によってソフトに破られる。3つの質量は次のようになる:
	\begin{equation}
		m_1 = A - B + \delta,\quad
		m_2 = A - B - \delta,\quad
		m_3 = A + 2B.
	\end{equation}
	スケーリング則 \(m \propto q^{2N}\) は、これらの質量を \(1\), \(r\), \(r^2\) に比例させる。ここで \(r = q^{N_\mu - N_e} = q^{39.5}\) である。\(A,B,\delta\) を消去すると、再び条件 \(Q = (1+r+r^2)/(1+\sqrt{r}+r)^2 = 2/3\) に導かれる。この方程式を満たす値 \(r \approx 0.2956\) は、まさに半整数指数 \(N_\mu - N_e = 39.5\) によって生成されたものである。
	
	したがって、小出公式は孤立した好奇心ではない; それは \(S_3\) 対称協調セルと普遍質量梯子の指紋である。核結合を支配する同じ位相シフト \(\sigma = 0.20\) (セクション~\ref{sec:sigma}) は、観測されたレプトン質量の理想的な \(Q = 2/3\) からの小さな逸脱も説明し、現象論的な円環を完全に閉じる。
	
	\subsection{ニュートリノセクターへの含意}
	
	シーソー機構または他の拡張がニュートリノ質量を同じ梯子に関連付けるならば、小出条件は3つの活性ニュートリノについての類似物を持つべきである。絶対ニュートリノ質量スケールに関する現在の限界は、すでに半整数梯子 (セクション~\ref{sec:formalism}) と整合する値を指し示している。ニュートリノ質量二乗差の高精度測定は、原理的に拡張された小出型関係をテストすることができ、協調パラダイムのさらなる相互検証を提供する。
	
	\section{議論}
	\label{sec:discussion}
	
	\subsection{協調深度の代数的構造: \(96 = 3 \times 16 \times 2\)}
	
	基準深度 \(L = 96\) は自然な因数分解を許容し、深い代数的階層を明らかにする:
	\[
	96 = 3 \times 16 \times 2.
	\]
	各因子は、YUCT フレームワーク内の異なる「協調ループ」に対応し、これは定数 \(\Sodd = 6/5\) と \(\Seven = 4/5\) を生成する代数的ループに類似している:
	
	\begin{itemize}
		\item \textbf{\(3\) — 世代。} 3つの世代は、大きな協調ループの3つの連続した巻きとして見ることができる。連続する世代間の質量比はほぼ \(3/2 = \Sodd/\Seven\) であり、世代間階層を基本ループ定数に直接結びつける。
		\item \textbf{\(16\) — \(SO(10)\) スピノル次元。} 単一世代内で、16個のワイルフェルミオン (右巻きニュートリノを含む) は閉じた代数的構造 — \(SO(10)\) スピノル表現 — を形成する。これは \emph{内部協調ループ} であり、その16の独立したチャネルが集合的に総深度に寄与する。
		\item \textbf{\(2\) — ヘリシティ / 粒子–反粒子。} この二項ループは CPT 対称性とスピン統計を反映する。それはスケーリング指数 \(\beta = 2 \times (1/3)\) の因子2として、そしてフェルミオン質量量子化 \(N_f\) の半整数ステップとして現れる。
	\end{itemize}
	
	したがって、総協調深度 \(L = 96\) は3つの入れ子になった代数的ループ (世代 (3), ゲージ構造 (16), スピン (2)) の次元の積である。この階層的因数分解は、同じ普遍比 \(3/2\) が弱-プランク階層 \(\sim (3/2)^{96}\) と世代間質量比の両方を支配する理由を説明する。また、YUCT が自然に \(SO(10)\) 大統一理論に埋め込まれることを強く示唆しており、16次元スピノル表現が基本協調多重項の役割を果たす。
	
	この代数的解釈は、数字96を単なる経験的入力から理論の構造的予測へと変え、時空の幾何学 (\(D=3\))、内部対称群、および量子場の統計を結びつける。
	
	\subsection{現代の階層問題との関連}
	
	LHC における階層問題の従来の解決策 (超対称性、複合性、または大きな余剰次元など) の証拠がないことは、自然さの広範な再評価を引き起こした。Koren が最近の包括的な論文 \cite{Koren2020} で述べているように、TeV スケールでの新しい可視状態の欠如は、問題を「Loerarchy 問題」に変えた: どのようにして赤外スケールが付随する構造なしに生成され得るのか？この危機は、低スケールの存在を近くの自由度の存在と結びつける有効場の理論 (EFT) の期待の明らかな頑健性によってさらに悪化している。
	
	YUCT フレームワークは、この EFT パラダイムからの急進的な脱却を提供する。二次発散を打ち消すために新しい対称性や粒子を呼び起こす代わりに、YUCT は質量という概念自体がグローバルな協調深度 \(L\) によって支配されると仮定する。弱いスケールは仮想的なパートナーを含む局所的な相殺によって安定化されるのではなく、むしろ標準模型のゲージ群と表現によって固定された数であるワイルフェルミオン自由度の総数 (\(L=96\)) によって \emph{決定} される。新しい TeV スケールのパートナーは不要であり、LHC のヌル結果に対する自然な説明を提供する。EFT の期待の明らかな違反は、したがって自然さの失敗としてではなく、普遍因子 \((3/2)^{96}\) を通じてプランクスケールを電弱スケールに直接結びつける、より深い非局所的な協調原理の証拠として再解釈される。
	
	\subsection{完全な協調データベースに向けて}
	
	ここで提示された結果は、すべての118元素に体系的に拡張することができる。すべての \(118\times 118\) ペアの完全な行列 \(C_{AB}\) は、半整数量子数 \(N_f\) とともに、材料発見と核物理学のための強力なツールを構成する。
	
	\subsection{限界と将来の作業}
	
	共鳴幅 \(\sigma = 0.20\) は経験的である; 基礎となる YUCT ダイナミクスからのその導出は未解決の問題である。熱化学データベース (例えば混合エンタルピー) に対する大規模な検証が必要である。分子や複雑な材料への適合性行列の拡張には、PLCM (周期的協調物質ラグランジアン) フレームワークを統合する必要がある。
	
	\subsection{プランク質量への逆スケーリング: 内部自己無撞着性チェック}
	\label{subsec:reverse}
	
	基本量子 $q = (3/2)^{1/3}$ を持つ経験的質量梯子は、自然な自己無撞着性テストを誘う: 電子質量 $m_e$ から始めてスケーリング因子 $q$ を繰り返し適用すると、プランク質量 $\Mpl = 1.2209\times10^{19}$ GeV に達するのに何ステップ $N_{\mathrm{Pl}}$ が必要か？真の基本スケールは、(小さな補正を除いて) 量子数の整数値または半整数値で到達するはずである。
	
	$m_e \cdot q^{N_{\mathrm{Pl}}} = \Mpl$ の定義を用いると、
	\begin{equation}
		N_{\mathrm{Pl}} = \frac{\ln(\Mpl/m_e)}{\ln q}
		= \frac{\ln\!\big(2.389\times10^{22}\big)}{\frac13\ln(3/2)}
		\approx 381.9.
		\label{eq:Npl}
	\end{equation}
	これは整数 $382$ に極めて近い。$N_{\mathrm{Pl}} = 382$ と設定すると、「予測された」プランク質量を得る:
	\begin{equation}
		\Mpl^{\mathrm{(ladder)}} = m_e \cdot q^{382}
		= m_e \cdot (3/2)^{382/3}
		\approx 2.61\times10^{22}\,m_e
		\approx 1.33\times10^{19}\,\text{GeV},
		\label{eq:Mpl_ladder}
	\end{equation}
	これは標準値からわずか $9\%$ しかずれていない。この外挿が乗法因子で382次、質量で22桁以上にわたることを考慮すると、整数への近接性は非常に非自明であり、スケーリング量子 $q$ がフィッティングの産物ではなく、自然の真の構造的特性であるという強力な内部証拠を提供する。
	
	さらに、整数 $382 = 2 \times 191$ は、完全な YUCT フレームワークで提案されている19次元幾何学との可能な接続を示唆している (19は親時空の基数; 因子2はヘリシティまたは粒子/反粒子の倍増に由来する可能性がある)。第一原理からの厳密な導出は未解決の問題のままであるが、この逆スケーリングテストは、経験的質量梯子が \emph{自己無撞着に} プランクスケールをその自然な上限境界として指し示していることを示している。
	
	\subsection{QCD からのハドロン質量創発との整合性}
	
	ジェファーソン研究所の最近の結果は、陽子質量の98\%以上が強い相互作用ダイナミクスから生じることを示している。YUCT では、これは協調深度 \(L\) の深化に対応する。陽子の有効深度 \(L_{\mathrm{eff}} \approx 108.5\) はそれを基本フェルミオンや原子核と同じ普遍質量スケール上に位置づけ、質量生成が協調原理によって普遍的に支配されることを示している。
	
	\subsection{なぜこの数学が量子力学の数学であるのか}
	\label{subsec:QM-as-YUCT}
	
	我々は今、最終的な概念的ギャップを埋める立場にある。すべての既知の粒子の質量を再現する代数的構造 (\(m = M_{\mathrm{Pl}}(2/3)^L\xi\) (有理数 $L$ と普遍的な $\beta=2/3$ を持つ)) は、量子力学の数学的形式の基礎となるのと正確に同じ代数的構造である。別個の「量子数学」は存在しない; ただ協調代数だけが存在し、それは何が参照スケールとして選ばれるかに応じて異なる観測量に投影される。
	
	\medskip
	\noindent\textbf{1. 量子化と離散性。}
	標準量子力学は、物理量が「量子化」されることを仮定する。YUCT では、離散性は自動的である:
	\[
	m_f = M_{\mathrm{Pl}} \left( \frac{2}{3} \right)^{L_f} \xi_f,
	\]
	ここで協調深度 $L_f$ は整数または半整数の値を取る。許容される質量は対数周期的な梯子を形成する。梯子が単一の基本比 $3/2 = S_{\mathrm{odd}}/S_{\mathrm{even}}$ によって生成されるという事実は、量子力学の量子化規則が、離散的で階層化された協調空間の対数射影であることを証明する。
	
	\medskip
	\noindent\textbf{2. 普遍指数と誤差法則。}
	普遍誤差法則 $\varepsilon \propto K_{\mathrm{eff}}^{-\beta}$ (with $\beta = 2/3$) は、伝送誤差の単なる経験的スケーリング関係ではない。d‑YPSDC の図では、量子場は協調場 $\Psi_{MN}$ である。この場の変動 — 「量子ゆらぎ」 — は同じ誤差法則によって支配される。したがって、量子力学の確率分布は、$K_{\mathrm{eff}}$ が有限である領域での辞書活性化誤差の統計的記述である。極限 $K_{\mathrm{eff}}\to\infty$ は古典的決定論を回復する; 有限の $K_{\mathrm{eff}}$ は量子理論の還元不可能な確率的特性を生成する。
	
	\medskip
	\noindent\textbf{3. プランク定数は協調スケールである。}
	基準深度 $L=96$ は弱スケール $m_W \sim M_{\mathrm{Pl}}(2/3)^{96}$ を固定する。$M_{\mathrm{Pl}}$ 自体が $\hbar$ を含むため、全質量スペクトルは正準交換関係の基礎となる同じ基本スケールに固定されている。YUCT から $\hbar$ を「導出」する必要はない; むしろ、$\hbar$ は無次元深度 $L$ と従来のエネルギー単位との間の変換因子であることが明らかにされる。それは人間の単位選択の歴史的産物であり、新しい基本定数ではない。
	
	\medskip
	\noindent\textbf{4. 互換性のない辞書エントリからの非可換性。}
	交換関係 $[x,p]=i\hbar$ は、相補的な性質を指す存在論的辞書の2つの異なるエントリを同時に活性化することの不可能性についての記述である。位置を調べる指標を送ることは、運動量レジスタを未指定のままにする。なぜなら、辞書は共役対に組織化されているからである (付録~X の量子–YPSDC 辞書参照)。この組織化を強制する代数的構造は、質量階層を組織化するのと同じ $3/2$ ループである。
	
	\medskip
	\noindent
	\textbf{結論。} YUCT の数学 — 代数的ループ、対数スケーリング、普遍指数 $\beta=2/3$ — は、協調優先存在論の視点から見た量子力学の数学である。追加の「量子」仮定は存在しない。シュレーディンガー方程式から正準交換関係に至るまで、すべての標準的な量子公式は、全体的な構造が代数的定数 $S_{\mathrm{odd}}=6/5$, $S_{\mathrm{even}}=4/5$, $\beta=2/3$ によって固定された階層的辞書内の特別目的の協調規則として得ることができる。量子論の「謎」は、量子化されているのがエネルギーや作用そのものではなく、粒子が存在する協調層の深さであることを認識するときに消える。
	
	\section{結論}
	\label{sec:conclusion}
	
	我々は、基本フェルミオン、ゲージボソン、軽いハドロン、および選択された重元素にまたがる質量と相互作用の統一的で協調に基づく記述を提示した。スケーリング量子 \(q = (3/2)^{1/3}\) の導入は自由パラメータの数を減らし、小さな補正を含めた後でフェルミオン質量に対してサブパーセントの精度を達成する。定量的共鳴相互作用係数 \(C_{AB}\) とそのヒートマップ可視化は、好ましい相互作用チャネルを明らかにする。この研究は、全元素の完全な協調データベースの基礎を築き、合理的材料設計のための YUCT の可能性を示す。
	
	\subsection{LHC データとの整合性と Loerarchy 逆説の解決}
	
	大型ハドロン衝突型加速器 (LHC) は、標準模型を超える物理学の徹底的な探索を実施したが、新しい粒子 (超対称パートナー、$Z'$ ボソン、複合状態、または大きな余剰次元) の確固たる信号は観測されていない。この TeV スケール構造の欠如は、「Loerarchy 問題」\cite{Koren2020} と呼ばれ、従来の自然さパラダイムに深刻な挑戦を提起する。超対称性、複合性、および余剰次元モデルでは、弱いスケールの安定性は、二次発散を打ち消すためにヒッグス質量付近の新しい状態を \emph{必要とする}。したがって、LHC のヌル結果はこれらのフレームワークの中心的な動機と直接矛盾する。
	
	階層問題に対する YUCT の解決策は根本的に異なる。弱いスケールは仮想的なパートナーを含む局所的な相殺によって安定化されるのではなく、むしろグローバルな協調深度 $L=96$ によって固定される。セクション~\ref{sec:formalism} で示したように、$L$ は自由パラメータではない: それは標準模型のワイルフェルミオン自由度の総数 (3世代 $\times 16$ $SO(10)$ スピノル成分 $\times 2$ ヘリシティ) に等しい。したがって、階層 $M_{\mathrm{Pl}}/m_W \sim (3/2)^{96}$ は、新しい物理学を待つ微調整問題ではなく、標準模型の場の内容の直接的な結果である。
	
	その結果、YUCT は TeV スケールのパートナーの不在を \emph{予測}する。LHC の「偉大な沈黙」は自然さの危機ではなく、YUCT パラダイムの直接的な経験的確証である。このフレームワークでは、例えば HL-LHC での超対称パートナーの発見は最小 YUCT 構造と調和させるのが難しいであろうが、従来のモデルは継続的な非観測によって反証される。
	
	さらに、弱いスケールを固定する同じ代数的構造が、すべての荷電フェルミオンの質量も支配する。半整数量子化規則 $m_f = m_e \cdot q^{N_f}$ (with $q = (3/2)^{1/3}$) は、9つの実験質量をパーセントレベルの精度と明確なパターンで再現する (表~\ref{tab:fermions_new} 参照)。この一致は12桁にわたり、自由パラメータの数を18から10に減らす。他のどのフレームワークも、ゲージ階層問題を同時に解決し、フェルミオンスペクトルを量子化し、LHC のヌル結果を自然に受け入れることはできない。
	
	\subsection{将来の実験的テスト}
	
	YUCT フレームワークは、現在および将来の実験データでテスト可能ないくつかの具体的で反証可能な予測を行う:
	
	\begin{itemize}
		\item \textbf{アップクォークとダウンクォークの質量:} 半整数量子化は、スケール $\mu \approx 2$ GeV で $m_u = 2.11$ MeV および $m_d = 4.75$ MeV を予測する。現在の格子 QCD 平均 \cite{PDG2024} ($m_u = 2.16 \pm 0.11$ MeV, $m_d = 4.67 \pm 0.09$ MeV) は合理的な一致を示している。将来の高精度格子計算はこれらの値の厳格なテストを提供するだろう。
		
		\item \textbf{HL-LHC での新しい粒子の不在:} YUCT は、(おそらくステライルニュートリノを除いて) High-Luminosity LHC や将来の100 TeV 衝突型加速器でも新しい基本粒子は発見されないと予測する。そのような粒子 (例えば $Z'$ ボソンや超対称パートナー) の発見は、最小 YUCT フレームワークを反証するであろう。
		
		\item \textbf{ニュートリノ質量の量子化:} スケーリング量子 $q = (3/2)^{1/3}$ はニュートリノセクターも支配すると予想される。絶対ニュートリノ質量スケールが一旦測定されれば (例えば KATRIN, Project 8, または宇宙論的サーベイによって)、質量は半整数梯子 $m_\nu = m_e \cdot q^{N_\nu}$ と整列するはずである。大きな逸脱は量子化規則の普遍性に挑戦するであろう。
		
		\item \textbf{重イオン衝突における軽い原子核の生成:} 適合性行列 $C_{AB}$ は、低い $C_{AB}$ を持つ原子核と比較して、高エネルギー重イオン衝突におけるアルファクラスター核 ($^4$He, $^{12}$C, $^{16}$O) の生成増加を予測する。LHC の ALICE 実験からの既存データはこの予測をテストするために分析できる。
	\end{itemize}
	
	YUCT フレームワークは、今後の実験によってテストされる具体的な予測を行う。特に、絶対ニュートリノ質量スケールは半整数量子梯子上にあると予測され、最も軽いニュートリノ質量は約 $0.023$ eV ($N_\nu = -125$) と予想される。現在の宇宙論的上限はこの値に近づいており、将来のサーベイ (DESI, CMB‑S4) と直接運動学的測定 (KATRIN++, Project 8) の組み合わせは、今後10年以内にこの予測を確認または反証するだろう。
	
	同様に、アップクォークとダウンクォークの質量に対する精密格子 QCD 計算は着実に改善されている。YUCT 予測 $m_u = 2.11$ MeV および $m_d = 4.75$ MeV ($\mu = 2$ GeV で) は現在の平均値と合理的に一致しており、格子の不確かさが減少するにつれて厳格にテストされるだろう。
	
	最後に、HL‑LHC での TeV スケールパートナーの不在は YUCT パラダイムをさらに裏付けるであろうが、例えば超対称粒子の発見はそれを強く不利にするであろう。これらは、反証不可能な数秘術から YUCT を区別する真の反証可能なテストである。
	
	\subsubsection{ニュートリノ質量スケール: 決定的テスト}
	\begin{table}[h]
		\centering
		\caption{YUCT 予測と LHC/PDG 測定の比較。}
		\normalsize
		\label{tab:lhc_comparison}
		\begin{tabular}{l c c c}
			\toprule
			粒子 & 実験 (GeV) & YUCT 予測 (GeV) & 一致 \\
			\midrule
			$W$ ボソン  & $80.377 \pm 0.012$ & $\approx 80$ & $<0.5\%$ \\
			$Z$ ボソン  & $91.1876 \pm 0.0021$ & $\approx 91.2$ & $<0.1\%$ \\
			ヒッグスボソン & $125.11 \pm 0.11$ & $\approx 124.2$ & $<0.7\%$ \\
			トップクォーク  & $172.69 \pm 0.30$ & $173.2$ & $<0.3\%$ \\
			ボトムクォーク & $4.18^{+0.03}_{-0.02}$ & $4.16$ & $<0.5\%$ \\
			ミューオン       & $0.105658$ & $0.1057$ & $<0.04\%$ \\
			\bottomrule
		\end{tabular}
	\end{table}
	
	YUCT フレームワークは、絶対ニュートリノ質量スケールについて具体的で反証可能な予測を行う: それはすべての荷電フェルミオンの質量を支配する同じ半整数量子梯子上に存在しなければならない。驚くべきことに、今日利用可能な最も厳しい実験的制約 — KATRIN による直接運動学的測定 (259日間のデータ) と DESI および ACT による精密宇宙論 — はこの予測と一致している。YUCT の自然な値 $m_\nu \approx 0.0234$ eV ($N_\nu = -125$) は現在の上限近くにあり、近い将来にテスト可能である。絶対ニュートリノ質量の将来の高精度測定は、YUCT をフェルミオン質量生成の正しい理論として確認するか、その反証をもたらすであろう — 真の科学的進歩の証である。
	
	\subsection*{現象論的フレームワークとしての YUCT の位置づけ}
	
	ヤクシェフ統一調整理論は、その現在の定式化では \textbf{現象論的フレームワーク} である。それは基本スケーリング量子 \(q = (3/2)^{1/3}\) を微視的作用や対称性原理から導出するのではなく、むしろ普遍指数 \(\beta = 2/3\) と代数的ループ定数 \(S_{\mathrm{odd}}=6/5\), \(S_{\mathrm{even}}=4/5\) に基づいてこの量子を \emph{仮定} する。第一原理からの導出 (例えば19次元協調幾何学から) は未解決の問題であり、進行中の研究の主題である。
	
	フレームワーク内のいくつかのパラメータ — 共鳴補正 \(\eta_f\)、ボソン因子 \(\xi\)、適合性幅 \(\sigma\) — は実験データにフィットされる。それらは現段階では理論によって予測されていない。自由パラメータの数は従来のフィットと比較して減少しているが、フレームワークはパラメータフリーではない。
	
	それにもかかわらず、ここで報告された中心的な経験的発見 — フェルミオン質量の半整数量子化、ワイル自由度の数 (\(L=96\)) からの弱いスケールの自然な創発、定量的適合性行列 \(C_{AB}\) — は頑健であり、これらのフィットされたパラメータの特定の値とは独立している。それらは、将来のいかなる基本理論も再現しなければならない質量スペクトルの中の一貫した代数的パターンを構成する。この点において、YUCT はバルマー公式やティティウス・ボーデの法則と類似の役割を果たす: それは自然の中の隠れた秩序を明らかにし、より深い基礎理論の探求を導く。
	
	\section*{謝辞}
	著者は、批判的議論を行った YUCT セミナーの参加者に感謝する。
	
	\section*{データの入手可能性}
	すべての表と計算は、YPSDC リポジトリ \url{https://github.com/Alexey-Yakushev-YUCT/YPSDC} で入手可能である。
	
	\begin{thebibliography}{99}
		\bibitem{AppendixY} A. V. Yakushev, ``Appendix Y: Mathematical Foundations of YUCT,'' in \emph{YUCT}, Zenodo (2026).
		\bibitem{AppendixJ} A. V. Yakushev, ``Appendix J: Complete Derivation of Standard Model Flavor Parameters from YUCT Topological Offsets,'' in \emph{YUCT}, Zenodo (2026).
		\bibitem{AppendixLambda} A. V. Yakushev, ``Appendix \(\Lambda\): Resolution of the Hierarchy and Cosmological Constant Problems within YUCT,'' in \emph{YUCT}, Zenodo (2026).
		\bibitem{AppendixFerra} A. V. Yakushev, ``Appendix Ferra1.2: Universal Coordination Constants for Ordered and Disordered Phases,'' in \emph{YUCT}, Zenodo (2026).
		\bibitem{PDG2024} S. Navas et al. (Particle Data Group), \emph{Review of Particle Physics}, Phys. Rev. D \textbf{110}, 030001 (2024).
		\bibitem{Massalski1990} T. B. Massalski (ed.), \emph{Binary Alloy Phase Diagrams}, 2nd ed., ASM International (1990).
		\bibitem{Okamoto1993} H. Okamoto, \emph{Phase Diagrams of Binary Magnesium Alloys}, ASM International (1993).
		\bibitem{Mokeev2025} V. I. Mokeev et al., ``Toward Understanding the Emergence of Hadron Mass,'' \emph{Symmetry} (Special Issue), 2025.
		\bibitem{Koren2020} S. Koren, \emph{The Hierarchy Problem: From the Fundamentals to the Frontiers}, Ph.D. dissertation, University of California, Santa Barbara (2020), arXiv:2009.11870 [hep-ph].
		\bibitem{CMS2024} CMS Collaboration, \emph{High-precision measurement of the W boson mass with the CMS experiment}, CERN (2024).
		\bibitem{Constantin2025} L. Constantin and F. Mahmoudi, \emph{The LHC has ruled out Supersymmetry -- really?}, Nucl. Phys. B \textbf{1018}, 117012 (2025), arXiv:2505.11251.
		\bibitem{YUCT} A. V. Yakushev, \emph{Yakushev Unified Coordination Theory (YUCT) V2.0}, Zenodo, DOI:10.5281/zenodo.18444599 (2026).
		\bibitem{KATRIN2025} KATRIN Collaboration, \emph{Direct neutrino-mass measurement based on 259 days of KATRIN data}, arXiv:2509.xxxxx (2025).
		\bibitem{DESI2025} DESI Collaboration, \emph{Constraints on Neutrino Physics from DESI DR2 BAO and DR1 Full Shape}, arXiv:2503.14744 (2025).
		\bibitem{T2KNOvA2025} T2K and NOvA Collaborations, \emph{Joint neutrino oscillation analysis from the T2K and NOvA experiments}, Nature \textbf{646}, 818-824 (2025).
		\bibitem{Feng2026} L. Feng et al., \emph{Measuring neutrino mass in light of ACT DR6 and DESI DR2}, arXiv:2603.10787 (2026).
		\bibitem{Parno2026} D. S. Parno (for the KATRIN Collaboration), \emph{Overview of Neutrino-Mass Determination}, arXiv:2601.01672 (2026).
		\bibitem{Koide1983} Y.~Koide, ``A New View of Quark and Lepton Mass Hierarchy,'' \emph{Phys. Rev. D} \textbf{28}, 252--254 (1983).
		\bibitem{Koide1982} Y.~Koide, ``Fermion‑Boson Two‑Body Model of Quark and Lepton Masses,'' \emph{Lett. Nuovo Cimento} \textbf{34}, 201--205 (1982).
		\bibitem{Foot1994} R.~Foot, ``A Note on the Koide Lepton Mass Relation,'' \emph{Phys. Lett. B} \textbf{326}, 307--311 (1994).
	\end{thebibliography}
	
\end{document}